\scanpage{88}

qui s'explicitent en termes de (78) en notant que tout drapeau de $J\cup L$ \inserted{intervenant dans (78)} s'écrit de façon unique
\[
(80)\qquad(j_{\unsure{1}}\ldots j_p\ l_1\ldots l_q)\qquad
\begin{gathered}
\unsure{i_0\ldots i_p|j_0\ldots j_p}\in J\\
l_1\ldots l_q\in L\\
p,q\geq1
\end{gathered}
\]
et \deleted{\unsure{ensuite}} on a des
\[(81)\]
\diagram{\begin{tikzcd}[ampersand replacement=\&,column sep=huge,row sep=large]
\Sigma_{j_1\ldots j_p\,l_1\ldots l_q}\arrow[r,hook]\arrow[d]\&\Sigma_{l_1\ldots l_q}\\
\Sigma_{j_1,\ldots,j_p}\&
\end{tikzcd}}

\uline{Remarque} Les $\Sigma_{J,L}$ ont un sens \inserted{(dans une catégorie isotopique convenable)} indépendamment des rétractions coniques locales autour des $X_i^*$, du moins en théorie topologique modérée, ainsi que les $\Sigma_{J,L}\hookrightarrow\Sigma_L$ --- et la description que j'en ai donné ici est adaptée à ce cas --- la structure $\Sigma_{J,L}\to\Sigma_J$ par contre utilise les \deleted{structures \unsure{de}\ \unsure{coniques}} \inserted{\unsure{fibrations}} locales. Quand on dispose $=$ des structures coniques locales (ce qui, en topologie modérée, et dans \unsure{la géom. algéb.|les géométries}, est une donnée supplémentaire ``anodine''), il est plus commode,
\margin{87}
